% Source: https://artofproblemsolving.com/community/c3h2472293_march_2021_problem_set
% Stable identity: aops-topic-2472293-q09
Label the roots $z_1,z_2,z_3,z_4$ so that $z_1z_2z_3z_4$ is the parallelogram in this order. The diagonals of a parallelogram bisect each other, so
\[z_1+z_3=z_2+z_4=t\]
for some $t$, and by Vieta's formulas $2t=z_1+z_2+z_3+z_4=a+16i$. Pairing the factors,
\[(x-z_1)(x-z_3)(x-z_2)(x-z_4)=\left(x^2-tx+u\right)\left(x^2-tx+v\right),\qquad u=z_1z_3,\ v=z_2z_4,\]
and expanding the right side gives $x^4-2tx^3+(t^2+u+v)x^2-t(u+v)x+uv$. Comparing coefficients:
\[t^2+u+v=147+306i,\qquad -t(u+v)=238-1610i,\qquad uv=-2000+bi.\]

\textbf{Finding $a$.} Eliminating $u+v$: $t^3-(147+306i)t-(238-1610i)=0$. Write $t=r+8i$ with $r=\dfrac a2$ real. Expanding, the real and imaginary parts are
\[r^3-339r+2210=0,\qquad 24r^2-306r-78=6(r-13)(4r+1)=0.\]
The second gives $r=13$ or $r=-\dfrac14$; only $r=13$ satisfies the first ($2197-4407+2210=0$, while for $r=-\dfrac14$ the left side is positive). So $t=13+8i$ and $\boxed{a=26}$. Then $t^2=105+208i$ and $u+v=42+98i$ (consistent with the $x$-coefficient, since $(13+8i)(42+98i)=-238+1610i$).

\textbf{Finding $b$.} So far $b$ is not determined; we use that the roots are Gaussian integers. Then $u,v$ are Gaussian integers with sum $42+98i$; write $u=(21+p)+(49+q)i$ and $v=(21-p)+(49-q)i$ with integers $p,q$. Since $b$ is real, $\mathrm{Re}(uv)=-2000$:
\[(441-p^2)-(2401-q^2)=-2000\iff p^2-q^2=40\iff(p-q)(p+q)=40.\]
The factors $p-q$ and $p+q$ have the same parity, so both are even, and $\{|p-q|,|p+q|\}=\{2,20\}$ or $\{4,10\}$. Hence $(|p|,|q|)=(11,9)$ or $(7,3)$. Swapping $u$ and $v$ changes the signs of both $p$ and $q$, so we may take $p>0$, leaving four cases.

The roots $z_1,z_3$ of $x^2-tx+u$ are $\dfrac{t\pm w}{2}$ with $w=z_1-z_3$ and $w^2=t^2-4u$, so $t^2-4u$ must be the square of a Gaussian integer, and then its absolute value is a perfect square integer. Similarly for $v$.
\begin{itemize}
\item $(p,q)=(11,9)$: $u=32+58i$, $t^2-4u=-23-24i$, with $|\,\cdot\,|^2=1105$, which is not a perfect square, so $|t^2-4u|$ is not an integer.
\item $(p,q)=(11,-9)$: $u=32+40i$, $t^2-4u=-23+48i$, with $|\,\cdot\,|^2=2833$, not a perfect square.
\item $(p,q)=(7,3)$: $u=28+52i$, $t^2-4u=-7$, which is not the square of a Gaussian integer ($(x+yi)^2=-7$ forces $xy=0$, and then $x^2=-7$ or $y^2=7$, both impossible).
\item $(p,q)=(7,-3)$: $u=28+46i$, $t^2-4u=-7+24i=(3+4i)^2$, so $z_1,z_3=\dfrac{13+8i\pm(3+4i)}2=8+6i,\ 5+2i$; and $v=14+52i$, $t^2-4v=49=7^2$, so $z_2,z_4=\dfrac{13+8i\pm7}{2}=10+4i,\ 3+4i$.
\end{itemize}
So the roots are $3+4i$, $5+2i$, $10+4i$, $8+6i$, a genuine parallelogram (sides $2-2i$ and $5+2i$), and
\[uv=(28+46i)(14+52i)=392-2392+(1456+644)i=-2000+2100i,\]
so $\boxed{b=2100}$.

\begin{asy}
usepackage("amsmath");
size(6cm);
defaultpen(fontsize(10pt));
import graph;
pair Z1=(3,4), Z2=(5,2), Z3=(10,4), Z4=(8,6), M=(6.5,4);
draw((-0.5,0)--(11.5,0),Arrow(6)); draw((0,-0.5)--(0,7.5),Arrow(6));
label("$\mathrm{Re}$",(11.5,0),E); label("$\mathrm{Im}$",(0,7.5),N);
for(int i=1;i<=7;++i) for(int j=1;j<=11;++j) dot((j,i),gray+linewidth(1.5));
fill(Z1--Z2--Z3--Z4--cycle,rgb(0.85,0.9,1));
draw(Z1--Z2--Z3--Z4--cycle,blue+linewidth(1));
draw(Z1--Z3,dashed); draw(Z2--Z4,dashed);
dot(Z1); dot(Z2); dot(Z3); dot(Z4); dot(M,red);
label("$3+4i$",Z1,W); label("$5+2i$",Z2,S); label("$10+4i$",Z3,E); label("$8+6i$",Z4,N);
label("$M$",M,dir(-65),red+fontsize(8pt));
label("$M=\dfrac{13}{2}+4i$",(6,-1.3),red+fontsize(9pt));
\end{asy}
